\scanpage{165}
% archive foliation: 164
\begin{center}- 34 -\end{center}
\hfill n$^\circ$ 297

Soient A un anneau topologique séparé, E un A-module topologique séparé. Soit $\unsure{\widetilde E|\widehat E}$ le groupe additif complété du groupe topologique commutatif E (\S\ 3, n$^\circ$5, th.2). L'application Z-bilinéaire $(\lambda,x)\longrightarrow\lambda x$ du produit $A\times E$ des groupes additifs A,E dans le groupe additif E se prolonge par continuité en une application Z-bilinéaire de $\widehat A\times\widehat E$ dans $\widehat E$ (n$^\circ$5, th.1), que nous désignerons encore par $(\lambda,x)\longrightarrow\lambda x$. En vertu du principe de prolongement des identités, on a encore $\lambda(\mu x)=(\lambda\mu)x$ pour $\lambda\in\widehat A$, $\mu\in\widehat A$, $x\in\widehat E$; la loi externe $(\lambda,x)\longrightarrow\lambda x$ définit donc sur $\widehat E$ une structure de $\widehat A$-module compatible avec la topologie de $\widehat E$; si en outre A admet un élément unité 1, cet élément est aussi élément unité de $\widehat A$, et si E est un A-module unitaire, on a encore $1.x=x$ pour $x\in\widehat E$ par le principe de prolongement des identités. Nous dirons que le $\widehat A$-module topologique $\widehat E$ ainsi défini est le \uline{complété} du A-module topologique E.

Soit E un A-module topologique, A et E n'étant pas nécessairement séparés; soient N (resp. F) l'adhérence de $\{0\}$ dans A (resp. E); N est un idéal bilatère de A (n$^\circ$3; prop.6) et F un sous-A-module de E (n$^\circ$1, prop.1); en outre on a par continuité $\lambda x\in F$ lorsque $\lambda\in N$ ou $x\in F$. On en déduit aussitôt, par passage aux quotients une application $(\mathring\lambda,\mathring x)\longrightarrow\mathring\lambda\mathring x$ de $(A/N)\times(E/F)$ dans $E/F$, et on vérifie aussitôt (grâce à la prop.24 du \S\ 2, n$^\circ$9) que cette application est loi de composition externe d'une structure de $(A/N$-module topologique sur $E/F$. Si l'on pose $B=A/N$, $L=E/F$, on dit que le B-module L est le module séparé \uline{associé} à E; son complété $\widehat L$ est un $\widehat B$-module topologique, qu'on appelle le module \uline{complété séparé} de E.

